\documentclass[10pt]{article}
\providecommand{\FIGDIR}{fig}
\newcommand{\DIGESTDATE}{31 July 2026}
% ===================== PM Research Watch - shared preamble =====================
\usepackage[T1]{fontenc}
\usepackage[utf8]{inputenc}
\usepackage{amsmath,amssymb}
\usepackage{tgpagella}
\usepackage{tgheros}
\renewcommand{\sfdefault}{qhv}
\renewcommand{\ttdefault}{lmtt}
\usepackage{microtype}
\usepackage[english]{babel}

\usepackage[a4paper,top=17mm,bottom=17mm,left=17mm,right=17mm,
            headsep=6mm,footskip=10mm]{geometry}
\usepackage{graphicx}
\usepackage{xcolor}
\usepackage{colortbl}
\usepackage{tikz}
\usetikzlibrary{calc}
\usepackage[most]{tcolorbox}
\usepackage{titlesec}
\usepackage{fancyhdr}
\usepackage{booktabs}
\usepackage{array}
\usepackage{longtable}
\usepackage{enumitem}
\usepackage{multicol}
\usepackage{ragged2e}
\usepackage[hidelinks]{hyperref}

% ------------------------------------------------------------------ palette
\definecolor{Ink}    {HTML}{14262E}
\definecolor{Deep}   {HTML}{0F4C5C}
\definecolor{Teal}   {HTML}{1F7A8C}
\definecolor{Sky}    {HTML}{4FA3C4}
\definecolor{Sage}   {HTML}{6FA287}
\definecolor{Amber}  {HTML}{E3A008}
\definecolor{Clay}   {HTML}{B5643C}
\definecolor{Coral}  {HTML}{D1495B}
\definecolor{Violet} {HTML}{6A5B9E}
\definecolor{Slate}  {HTML}{4A6572}
\definecolor{Paper}  {HTML}{FBF9F5}
\definecolor{Mist}   {HTML}{EEEAE2}
\definecolor{Rule}   {HTML}{D9D3C9}

\pagecolor{Paper}
\color{Ink}

\hypersetup{colorlinks=true, linkcolor=Deep, urlcolor=Teal, citecolor=Teal,
            pdftitle={PM Research Watch}, pdfauthor={Automated daily digest}}

% ------------------------------------------------------------------ headings
\titleformat{\section}
  {\sffamily\bfseries\large\color{Deep}}{}{0pt}
  {}[{\vspace{1.2mm}\color{Amber}\titlerule[1.1pt]}]
\titlespacing*{\section}{0pt}{5.5mm}{2.8mm}

\titleformat{\subsection}
  {\sffamily\bfseries\normalsize\color{Teal}}{}{0pt}{}
\titlespacing*{\subsection}{0pt}{3.6mm}{1.4mm}

\setlist[itemize]{leftmargin=4.2mm, itemsep=1.1mm, topsep=1.3mm, parsep=0pt,
                  label=\textcolor{Teal}{\footnotesize$\blacktriangleright$}}

% ------------------------------------------------------------------ header/footer
\pagestyle{fancy}
\fancyhf{}
\renewcommand{\headrulewidth}{0pt}
\renewcommand{\footrulewidth}{0pt}
\fancyhead[L]{\sffamily\scriptsize\color{Slate}PM RESEARCH WATCH\;\textcolor{Rule}{\textbar}\;\DIGESTDATE}
\fancyhead[R]{\sffamily\scriptsize\color{Slate}Particulate matter monitoring \& health effects}
\fancyfoot[C]{\sffamily\scriptsize\color{Slate}\thepage}

% ------------------------------------------------------------------ components
% metric tile
\newcommand{\metric}[3]{%
  \begin{tcolorbox}[enhanced, colback=white, colframe=Rule, boxrule=0.5pt,
    arc=1.6mm, left=2.2mm, right=2.2mm, top=1.8mm, bottom=1.8mm,
    borderline west={1.7mm}{0pt}{#3}, nobeforeafter, valign=center,
    height=17mm, width=\linewidth]
    \RaggedRight
    {\sffamily\bfseries\fontsize{17}{17}\selectfont\color{#3}#1}\\[0.5mm]
    {\sffamily\fontsize{6.5}{7.4}\selectfont\color{Slate}#2}
  \end{tcolorbox}}

% pull-quote / key insight box
\newtcolorbox{keybox}[1][]{enhanced, colback=Mist, colframe=Mist,
  arc=1.6mm, left=3.2mm, right=3.2mm, top=2.4mm, bottom=2.4mm,
  borderline west={1.7mm}{0pt}{Amber}, #1}

% subtopic band
\newcommand{\band}[2]{%
  \vspace{3.2mm}\noindent
  \begin{tcolorbox}[enhanced, colback=#2, colframe=#2, arc=1.2mm,
    left=2.6mm, right=2.6mm, top=1.1mm, bottom=1.1mm, boxrule=0pt,
    width=\linewidth, nobeforeafter]
    {\sffamily\bfseries\small\color{white}#1}
  \end{tcolorbox}\vspace{1.6mm}}

% one paper entry: {short cite}{title}{journal}{doi}{body}
\newcommand{\paperentry}[5]{%
  \par\vspace{1.9mm}\noindent
  {\sffamily\bfseries\footnotesize\color{Ink}#1}\;%
  {\sffamily\fontsize{7}{8}\selectfont\color{Slate}#3}\\[0.4mm]
  {\footnotesize\itshape\color{Deep}#2}\\[0.6mm]
  {\footnotesize #4}\\[0.3mm]
  {\fontsize{6.6}{7.6}\selectfont\color{Slate}\texttt{doi:}\,\href{https://doi.org/#5}{\texttt{\detokenize{#5}}}}
  \par}

% caption style
\newcommand{\figcap}[1]{\par\vspace{1.0mm}
  {\fontsize{7.2}{8.6}\selectfont\color{Slate}\textsf{\textbf{Fig.}}\; #1\par}\vspace{1.5mm}}

\providecommand{\permil}{\textperthousand}

\setlength{\parindent}{0pt}
\setlength{\parskip}{0pt}
\renewcommand{\arraystretch}{1.18}


\begin{document}

% ==================================================================== MASTHEAD
\thispagestyle{empty}
\begin{tikzpicture}[remember picture, overlay]
  \fill[Deep] (current page.north west) rectangle
      ($(current page.north east)+(0,-67mm)$);
  \fill[Amber] ($(current page.north west)+(0,-67mm)$) rectangle
      ($(current page.north east)+(0,-68.4mm)$);
  \foreach \x/\y/\r in {22/12/0.9, 41/26/1.5, 63/9/0.7, 88/31/1.1, 112/17/1.9,
                        138/29/0.8, 159/13/1.3, 178/54/1.0, 196/21/1.6,
                        30/44/0.6, 74/57/1.2, 125/48/0.9, 168/6/0.8, 191/59/0.7,
                        52/18/1.0, 100/60/0.6, 148/20/1.1, 66/49/1.4, 205/45/0.8}
    \fill[Sky, opacity=0.30] ($(current page.north west)+(\x mm,-\y mm)$) circle (\r mm);
\end{tikzpicture}

\vspace*{4mm}
{\sffamily\fontsize{7.6}{9}\selectfont\color{Sky}%
 AUTOMATED DAILY LITERATURE DIGEST\;\textbullet\;PhD RESEARCH WATCH\par}
\vspace{2.2mm}
{\sffamily\bfseries\fontsize{27}{29}\selectfont\color{white}%
 Particulate Matter\par}
{\sffamily\fontsize{27}{29}\selectfont\color{Sky}%
 Monitoring \& Health Effects\par}
\vspace{5.0mm}
{\sffamily\fontsize{8.4}{10}\selectfont\color{white}%
 Issue 2026-07-31\;\;\textcolor{Amber}{\textbar}\;\;Entry window 31 July 2026\;\;%
 \textcolor{Amber}{\textbar}\;\;11 records in scope\par}
\vspace{18mm}

% -------------------------------------------------------------- metric strip
\noindent
\begin{minipage}[t]{0.190\linewidth}\metric{11}{RECORDS IN SCOPE\\(7 SUBTOPICS)}{Deep}\end{minipage}\hfill
\begin{minipage}[t]{0.190\linewidth}\metric{8}{EFFECT ESTIMATES\\EXTRACTED}{Teal}\end{minipage}\hfill
\begin{minipage}[t]{0.190\linewidth}\metric{4}{RECORDS FROM CHINA\\(36\% OF CORPUS)}{Sky}\end{minipage}\hfill
\begin{minipage}[t]{0.190\linewidth}\metric{2}{EXPOSURE ASSESSMENT\\\& MODELLING}{Amber}\end{minipage}\hfill
\begin{minipage}[t]{0.190\linewidth}\metric{0}{SENSING /\\INSTRUMENTATION}{Clay}\end{minipage}

\vspace{6mm}

% -------------------------------------------------------------- signal of day
\section*{\sffamily\bfseries\color{Deep}Signal of the day}
\vspace{-1.5mm}{\color{Amber}\rule{\linewidth}{1.1pt}}\vspace{2.5mm}

\begin{keybox}
{\sffamily\bfseries\small\color{Deep}Road-network proxies fail against personal PM$_{2.5}$
in sub-Saharan Africa --- and they fail with the wrong sign, not merely weakly.}\\[1.4mm]
\footnotesize
Nyoni et al.\ (\textit{Air Qual Atmos Health}) validated three of the exposure proxies
that epidemiology routinely substitutes for measurement --- population-weighted road
network density, Euclidean distance to highway, Euclidean distance to main road ---
against \textbf{a year of filter-based personal PM$_{2.5}$ from 343 postpartum
participants} in The Gambia, Kenya and Mozambique. Weighted road density correlated with
personal PM$_{2.5}$ only in Mozambique ($\rho = 0.351$, $p = 0.005$); Kenya and The Gambia
were null ($\rho = -0.041$ and $-0.020$). Distance-to-highway was \emph{positive} in
The Gambia and Mozambique but \textbf{significantly negative in Kenya}
($\rho = -0.224$, $p = 0.018$) --- the proxy inverts across countries. Single-variable
random forests never exceeded $R^2 = 0.45$; combining all three reached
$R^2 = 0.52$ (RMSE 31.7\,$\mu$g/m$^3$) in Kenya and 0.60 (8.9\,$\mu$g/m$^3$) in
Mozambique, with a global model at 0.46. An RMSE of 30\,$\mu$g/m$^3$ is larger than the
entire annual mean in most European cohorts. The operational reading for
distributed sensing is direct: in settings where household energy and biomass burning
dominate the personal exposure budget, traffic geometry is not a transferable
surrogate, and the case for physically deploying low-cost nodes rather than
regressing on road maps is strongest exactly where reference infrastructure is
thinnest.
\end{keybox}

\vspace{3mm}

% -------------------------------------------------------------- exec brief
\section*{\sffamily\bfseries\color{Deep}What changed today}
\vspace{-1.5mm}{\color{Amber}\rule{\linewidth}{1.1pt}}\vspace{2.5mm}

\footnotesize
\begin{itemize}
\item \textbf{Two independent records say the exposure variable is the problem, not the
model.} Nyoni et al.\ show road-geometry proxies inverting in sign against measured
personal PM$_{2.5}$ between African countries; Knapova et al.\ show that step counts fall
361 per 1-SD PM$_{10}$ rise, which makes time-activity patterns \emph{endogenous} to
concentration. Taken together they bracket the same failure from both ends: the spatial
surrogate is wrong, and the mobility assumption that would let you correct it is wrong
too.

\item \textbf{PM mass was the null term in its own model, twice.} Zhu et al.\ report
PM$_{2.5}$ RR 1.004 (0.991--1.017) and PM$_{10}$ 1.004 (0.999--1.010) per
10\,$\mu$g/m$^3$ for pneumonia admissions against NO$_2$ at 1.221 (1.174--1.271); Yue
et al.\ report no particulate term at all and attribute childhood bronchopneumonia to
NO$_2$ and SO$_2$. In both cases the pollutants share a combustion source and a seasonal
cycle and no multi-pollutant model is fitted, so these are statements about correlation
structure rather than about which pollutant acts --- but the operational consequence
stands: a mass-only monitor would have recorded nothing on either study's strongest days.

\item \textbf{Mixture methods are outrunning the sample sizes they are applied to.}
Chen et al.\ run quantile g-computation over 29 hazards in 7{,}092 adults and recover
asthma odds of 1.35 (1.06--1.73) per quintile \textbf{only} in high-social-vulnerability
tracts; Sezer et al.\ run a six-pollutant WQS index on 127 cleft cases and report ozone
dominance with no bootstrap weight distribution; Zhou et al.\ run WQS on urinary
metabolites measured \emph{at or after} diagnosis. The technique is the same; the
credibility is not, and it tracks sample size and exposure timing rather than the model.

\item \textbf{Zero sensing or instrumentation records again.} This is a harvest defect,
not a real absence: the July monthly rollup recovered 23 sensing records for the month
against 11 across all five daily issues combined. The pipeline still has no date-windowed
channel into \textit{Atmos.\ Meas.\ Tech.} or \textit{Atmos.\ Chem.\ Phys.} See the
provenance box.
\end{itemize}

\clearpage

% ==================================================================== FIGURES
\section*{\sffamily\bfseries\color{Deep}Corpus analytics}
\vspace{-1.5mm}{\color{Amber}\rule{\linewidth}{1.1pt}}\vspace{2.5mm}

\noindent\includegraphics[width=\linewidth]{\FIGDIR/f1_subtopics.png}
\figcap{Subtopic assignment of the 11 in-scope records, single-label by dominant
contribution. Seven of the ten subtopics are populated; Neuro / mental health,
Cardiovascular \& metabolic and Sensing, forecasting \& instrumentation are empty this
window. Colours are fixed across issues so a subtopic keeps its hue in the rollups.}

\vspace{3mm}

\noindent\includegraphics[width=\linewidth]{\FIGDIR/f2_design_endpoint.png}
\figcap{Left --- study architecture. Observational designs account for 6/11, split 3
cross-sectional (two case-control, one exposome-wide) and 3 acute time-series or panel.
Two measurement campaigns, one modelling record, one evidence synthesis and one
experimental-toxicology record complete the set. Right --- health endpoint. Two records
carry no human health outcome: the exposure-proxy validation (plotted as \emph{No health
endpoint}) and the indoor-air characterisation (plotted as \emph{Microenvironment}).}

\vspace{3mm}

\noindent\includegraphics[width=\linewidth]{\FIGDIR/f3_metric_geography.png}
\figcap{Left --- particle metric. PM$_{2.5}$ alone in 4 records, PM$_{2.5}$ with
PM$_{10}$ in 3, PM$_{10}$ alone in 1, and 3 records use an unspeciated or non-ambient
particulate metric (diesel SRM 1650b, indoor construction PM, traffic gases standing in
for PM). Right --- geography. China supplies 4 of 11 records; the single sub-Saharan
African record is the exposure-proxy validation study.}

\vspace{3mm}

\noindent\includegraphics[width=\linewidth]{\FIGDIR/f6_lifecourse.png}
\figcap{Life-course windows addressed. Counts are non-exclusive. Working age leads at 3;
no record addresses adolescence. Records with no identifiable human life-course window
--- the ex-vivo cardiac preparation and the rodent meta-analysis --- are excluded from
the tally.}

\clearpage

\noindent\includegraphics[width=\linewidth]{\FIGDIR/f4_heatmap.png}
\figcap{Subtopic $\times$ study-architecture cross-tabulation. With 11 records the matrix
is sparse and little should be read from individual cells; the one structural feature
worth noting is the empty Sensing row, which is a harvest artefact rather than a property
of the literature.}

\vspace{3mm}

\noindent\includegraphics[width=\linewidth]{\FIGDIR/f5_forest.png}
\figcap{The eight quantitative effect estimates in the corpus on a common ratio axis.
\textbf{The exposure contrasts are not commensurable} and the plot must not be read as a
meta-analysis: increments range from per 5\,$\mu$g/m$^3$ to per 10\,$\mu$g/m$^3$ to
quartile and quintile contrasts to a WQS index, and the outcome scale mixes RR and OR.
Note that the two PM-specific estimates (Zhu et al., PM$_{2.5}$ and PM$_{10}$ pneumonia
admissions) are the smallest in the set and are null or borderline, while every large
estimate belongs to a gaseous co-pollutant or a mixture index.}

\clearpage

% ==================================================================== DIGEST
\section*{\sffamily\bfseries\color{Deep}Paper-level digest}
\vspace{-1.5mm}{\color{Amber}\rule{\linewidth}{1.1pt}}\vspace{2.5mm}

\band{Exposure assessment \& modelling\hfill 2 records}{Amber}

\paperentry{Nyoni et al.\ 2026}
{Performance of road-traffic-based exposure proxies against personal PM$_{2.5}$
measurements in three sub-Saharan African countries}
{Air Qual Atmos Health}
{One year of filter-based outdoor personal PM$_{2.5}$ from 343 postpartum participants in
The Gambia, Kenya and Mozambique, used to validate population-weighted road network
density (wRND), Euclidean distance to highway (EH) and to main road (EM). Spearman
correlations were weak and \textbf{sign-inconsistent across countries}: wRND significant
only in Mozambique ($\rho = 0.351$), EH positive in The Gambia (0.335) and Mozambique
(0.292) but \textbf{negative in Kenya ($-0.224$, $p = 0.018$)}. Single-proxy random
forests gave $R^2 < 0.45$ everywhere; combined models reached $R^2 = 0.52$
(RMSE 31.7\,$\mu$g/m$^3$) in Kenya and 0.60 (8.9\,$\mu$g/m$^3$) in Mozambique. The
sampled population --- postpartum women --- is not representative of general mobility
patterns, and only outdoor-filtered personal exposure was measured, so the household
cooking contribution that most plausibly drives the proxy failure is inferred rather than
partitioned. Still the strongest available empirical case that road-geometry proxies
require local validation before use, and a direct argument for node deployment over GIS
surrogacy.}
{10.1007/s11869-026-02060-y}

\paperentry{Chen et al.\ 2026}
{Geospatial-based neighborhood hazards and asthma in adults: interaction with social
vulnerability}
{J Hazard Mater}
{Exposome-wide association study over 29 residential hazards --- air pollutants,
industrial VOC emissions, highway density and proximity to 13 pollution-source classes
--- in 7,092 North Carolina adults (2013--2022), stratified at the median census-tract
Social Vulnerability Index, with quantile g-computation for the mixture. In the overall
population only proximity to dry-cleaning-solvent-contaminated sites was associated with
asthma. In high-SVI tracts, \textbf{eight hazards --- predominantly PM$_{2.5}$
constituents and VOCs --- were positively associated at FDR $<$ 0.05}, and each
one-quintile increase in the hazard mixture carried \textbf{35\% higher asthma odds
(95\% CI 6--73\%)}. Toluene, dry-cleaning proximity, PM organic matter and carbon
monoxide each contributed over 10\% of the joint association. Hazard correlations were
themselves stronger in high-SVI areas, which is both the substantive finding and the
analytical problem: with collinearity differing by stratum, the g-computation weights are
not directly comparable across strata. Cross-sectional, self-reported
physician-diagnosed asthma, residence-only exposure.}
{10.1016/j.jhazmat.2026.143083}

\band{Respiratory \& allergic\hfill 2 records}{Sage}

\paperentry{Yue et al.\ 2026}
{Short-term traffic-related gaseous pollutant exposure and childhood bronchopneumonia
hospitalization in Guangzhou}
{Front Public Health}
{20,623 paediatric bronchopneumonia admissions (ICD-10 J18, ages 0--14) at a single
tertiary hospital in Huadu District, Guangzhou, 2014--2023, against daily NO$_2$ and
SO$_2$ in a quasi-Poisson generalised additive model. Per 5\,$\mu$g/m$^3$, peak
single-day excess risk at lag 4 was \textbf{2.74\% (95\% CI 1.73--3.75) for NO$_2$} and
\textbf{8.81\% (4.94--12.82) for SO$_2$}; over the 0--7-day moving average these rose to
\textbf{6.83\% (4.93--8.77)} and \textbf{17.21\% (10.18--24.70)}. Children aged 6--14 and
boys were consistently more susceptible. Two problems: no particulate term is reported at
all, so the gaseous estimates absorb whatever correlated PM signal exists and cannot be
read as pollutant-specific; and a single-hospital catchment in a ten-million-person
megacity leaves the denominator undefined and exposed to referral-pattern drift over a
decade. Included because the lag structure and the age--sex susceptibility gradient are
informative, not because the point estimates are attributable.}
{10.3389/fpubh.2026.1885590}

\paperentry{Zhu et al.\ 2026}
{Air pollution and pneumonia hospitalization in Qingyang, China: a time-series analysis}
{Medicine (Baltimore)}
{Five years of pneumonia admissions (2015--2019) in an inland Chinese energy city against
six pollutants, generalised additive plus distributed-lag nonlinear models. Per
10\,$\mu$g/m$^3$: \textbf{PM$_{2.5}$ RR 1.004 (0.991--1.017)} and \textbf{PM$_{10}$ 1.004
(0.999--1.010)} --- both null --- against \textbf{NO$_2$ 1.221 (1.174--1.271)},
\textbf{SO$_2$ 1.035 (1.022--1.048)}, and CO 1.382 (1.230--1.552) per 1\,mg/m$^3$.
Effects peaked at lag 1 for PM$_{2.5}$ and lag 0--7 for the gases; men were more
vulnerable to SO$_2$, women to CO, and children $\leq$14 formed the high-risk group. The
single-pollutant framing is the limitation --- in a heating-dominated city PM$_{2.5}$,
SO$_2$, NO$_2$ and CO share a common combustion source and a common seasonal cycle, and
no multi-pollutant model is presented, so the apparent NO$_2$ dominance is a statement
about correlation structure rather than about which pollutant acts. Useful mainly as a
demonstration that PM mass can be the null term in its own time-series.}
{10.1097/MD.0000000000049962}

\band{Mechanistic toxicology\hfill 2 records}{Violet}

\paperentry{Balu et al.\ 2026}
{SRM 1650b administration to isolated rat heart aggravates ischemia-reperfusion injury
via mitochondrial dysfunction}
{J Biochem Mol Toxicol}
{Langendorff-perfused male Wistar hearts exposed to NIST SRM 1650b heavy-duty diesel
particulate at 10, 100 and 300\,$\mu$g/mL before 30\,min ischemia and 60\,min
reperfusion. Dose-dependent deterioration of haemodynamics and elevated tissue injury,
with increased oxidative stress in tissue and isolated mitochondria, reduced bioenergetic
enzyme activity and respiratory efficiency, downregulated \textit{Pgc-1$\alpha$},
\textit{Tfam}, \textit{Polg}, \textit{Fis1}, \textit{Mfn1} and \textit{Pink1}, lower mtDNA
copy number, and suppressed PI3K/Akt signalling that fell further after the ischemic
challenge. The design isolates a direct cardiomyocyte effect cleanly --- no autonomic,
pulmonary or systemic inflammatory route is available in an ex-vivo preparation --- which
is both its strength and its ceiling. The concentrations are far above any inhaled dose
delivered to myocardium in vivo, and direct coronary perfusion bypasses the alveolar
barrier entirely, so this establishes a mechanism that \emph{can} operate, not one that
does at ambient exposure.}
{10.1002/jbt.71038}

\paperentry{Pan et al.\ 2026}
{Effects of air pollution exposure on nasal mucosal immune-inflammatory markers in animal
models of allergic rhinitis}
{Front Pharmacol}
{PRISMA systematic review and random-effects meta-analysis of 18 rodent studies,
standardised mean differences for nasal mucosal markers. Air pollutant exposure was
associated with increased eosinophils and elevated IL-4, IL-5, IL-13 and OVA-specific IgE
--- a coherent enhancement of type 2 inflammation --- with scattered evidence of
epithelial barrier loss (reduced ZO-1) and innate activation (IL-1$\beta$, NLRP3).
IFN-$\gamma$ showed no consistent direction. Heterogeneity was substantial and the
authors decline to claim a causal pathway. The pooling is the weak point: 18 studies
spanning diesel exhaust particulate, PM$_{2.5}$, PM$_{10}$ and mixed urban aerosol at
uncontrolled doses and durations, in an ovalbumin-sensitisation model whose relationship
to human allergic rhinitis is itself contested, cannot yield a transferable effect size.
No dose--response meta-regression is attempted, which is the analysis that would have
justified the synthesis. Value here is as a marker map for the type 2 axis, not as an
effect estimate.}
{10.3389/fphar.2026.1870023}

\band{Other clinical endpoints\hfill 2 records}{Slate}

\paperentry{Zhou et al.\ 2026}
{Environmental exposome study for prostate cancer in elderly men from Western China}
{BMJ Open}
{580 histologically confirmed localised prostate cancers after radical prostatectomy,
matched 1:2 to 1,160 cancer-free controls from the WeEndPd cohort, with urinary
bisphenols and phthalate metabolites by LC-MS/MS and geocoded five-year SO$_2$ and
PM$_{2.5}$. Highest-quartile five-year SO$_2$ (9.37--28.79\,$\mu$g/m$^3$) gave
\textbf{OR 1.65 (95\% CI 1.08--2.72)} against the lowest; MECPP, MEHHP, MEHP, BPA and BPZ
were all higher in cases; the WQS mixture of endocrine-disrupting chemicals gave
\textbf{OR 1.26 (1.20--1.33)}. The PM$_{2.5}$ term is reported but is not the headline.
Two structural problems: urinary EDC metabolites have half-lives of hours to days and
were measured at or after diagnosis, so they cannot index the aetiologically relevant
window --- this is reverse-causation-prone by construction; and cases are surgical
patients from one tertiary centre, which selects on operability and on health-seeking
behaviour that also tracks residential air quality.}
{10.1136/bmjopen-2025-110957}

\paperentry{Knapova et al.\ 2026}
{Daily relationship between air pollution and physical activity in an industrial region}
{Appl Psychol Health Well Being}
{693 adults in the industrial Ostrava region of the Czech Republic, 9,702 person-days,
device-measured step counts against daily PM$_{10}$ over 14 days, multilevel Bayesian
models with within-person centring and adjustment for temperature, sunshine and
precipitation. On days when PM$_{10}$ exceeded a participant's own 14-day mean, step
counts fell ($b = -0.027$, 95\% CI $-0.047$ to $-0.006$); \textbf{a 1-SD PM$_{10}$
increase corresponded to 361 fewer steps}. Self-reported tendency to monitor air quality
did not moderate the effect, which argues the response is not deliberate avoidance
behaviour. Exposure is ambient monitor data assigned to residence, not personal, so the
estimate confounds pollution with unmeasured weather and visibility cues; 14 days is
short for a within-person design; and non-probabilistic quota sampling limits transport.
The finding matters for exposure modelling regardless: \textbf{behavioural response makes
time-activity patterns endogenous to concentration}, which biases any exposure model that
treats mobility as exogenous.}
{10.1111/aphw.70186}

\band{Burden, policy \& mitigation\hfill 1 record}{Coral}

\paperentry{Qian et al.\ 2026}
{Stronger actions are needed to advance PM$_{2.5}$ mitigation and its health benefits
under population aging in China}
{J Hazard Mater}
{A new daily 1\,km PM$_{2.5}$ product for mainland China 2015--2024, used to track
concentration, WHO AQG exceedance and population-weighted exposure, with attributable
cardiovascular mortality and a decomposition isolating the ageing offset. National
PM$_{2.5}$ fell \textbf{35.57 $\rightarrow$ 20.62\,$\mu$g/m$^3$} (population-weighted
$-37.35\%$), and the area above the WHO interim target of 35\,$\mu$g/m$^3$ collapsed from
51.07\% to 1.49\% --- yet \textbf{over 99\% of the population still lives above the
5\,$\mu$g/m$^3$ guideline}. Attributable CVD deaths fell from 2.16 (1.64--2.58) million
in 2015 to 1.55 (1.14--1.91) million in 2022, then \textbf{rebounded to 1.76 (1.30--2.15)
million in 2024}. The share of health benefit cancelled by population ageing rose from
around 40\% before 2020 to \textbf{56.54\% in 2024}, exceeding 70\% in some subregions.
The burden estimates inherit exposure--response functions from external cohorts, so the
uncertainty intervals understate structural uncertainty, and the 1\,km product's own
error is not propagated into the mortality intervals. The demographic-offset framing is
the contribution and it generalises beyond China.}
{10.1016/j.jhazmat.2026.143101}

\band{Reproductive \& developmental\hfill 1 record}{Sky}

\paperentry{Sezer et al.\ 2026}
{Preconception and early pregnancy exposure to ambient air pollution and risk of cleft
lip and/or palate}
{Birth Defects Res}
{Hospital-based case-control study, 127 nonsyndromic CLP cases against 271 controls, with
PM$_{2.5}$, PM$_{10}$, O$_3$, NO$_x$, SO$_2$ and CO assigned for the three months before
conception and the first trimester. Higher O$_3$, PM$_{2.5}$ and PM$_{10}$ were each
associated with elevated CLP risk in single-pollutant models; \textbf{O$_3$ was the
dominant contributor in the WQS mixture} and the only exposure stable across windows and
model specifications, while traffic-related pollutants did not survive multi-pollutant
adjustment. The study is underpowered for a mixture analysis --- 127 cases against a
six-pollutant WQS index invites weight instability, and no bootstrap weight distribution
is reported --- and hospital-based control selection risks referral bias. Read as
hypothesis-generating. The methodological point that carries: with correlated pollutants
and a small case series, the pollutant that wins the mixture is often the one with the
widest exposure contrast, not the one with the strongest biology.}
{10.1002/bdr2.70102}

\band{Occupational \& indoor exposure\hfill 1 record}{Clay}

\paperentry{Park et al.\ 2026}
{Environmental evaluation of wood-based remodeling in aging buildings: indoor air quality
and life-cycle carbon}
{Sci Total Environ}
{Mock-up chamber characterisation of softwood finishes plus field monitoring in an
occupied childcare facility and an unoccupied university classroom --- before
remodelling, 24\,h after, and at four weeks --- covering formaldehyde, anthropogenic and
biogenic VOCs, particulate matter and total airborne bacteria, with a life-cycle global
warming potential assessment. $\alpha$-Pinene dominated the biogenic emission profile.
In the occupied, ventilated building PM and bacterial counts fell or stabilised and most
VOCs declined within four weeks; in the unoccupied, poorly ventilated classroom
\textbf{$\alpha$-pinene rose from 28.7 to 1713.6\,$\mu$g/m$^3$ over four weeks} --- a
60-fold accumulation. Net GWP in that building fell to 37.7\,kgCO$_2$eq, 9--11 times
below non-wood alternatives, on biogenic carbon storage. With $n = 2$ buildings differing
simultaneously in occupancy, ventilation and use, those factors cannot be separated, and
the PM measurements are reported only qualitatively. The transferable point is the
ventilation dependence: a material acceptable in one air-exchange regime is not in
another, which argues for continuous indoor monitoring rather than commissioning-time
spot checks.}
{10.1016/j.scitotenv.2026.182065}

\clearpage

% ==================================================================== TABLE
\section*{\sffamily\bfseries\color{Deep}Corpus register}
\vspace{-1.5mm}{\color{Amber}\rule{\linewidth}{1.1pt}}\vspace{2.5mm}

{\fontsize{6.7}{8.4}\selectfont
\setlength{\tabcolsep}{2.4pt}
\begin{longtable}{@{}>{\RaggedRight}p{25mm}>{\RaggedRight}p{41mm}>{\RaggedRight}p{27mm}%
>{\RaggedRight}p{23mm}>{\RaggedRight}p{18mm}>{\RaggedRight}p{15mm}@{}}
\toprule
\textsf{\textbf{First author}} & \textsf{\textbf{Focus}} & \textsf{\textbf{Journal}} &
\textsf{\textbf{Design}} & \textsf{\textbf{Metric}} & \textsf{\textbf{Region}}\\
\midrule
\endfirsthead
\toprule
\textsf{\textbf{First author}} & \textsf{\textbf{Focus}} & \textsf{\textbf{Journal}} &
\textsf{\textbf{Design}} & \textsf{\textbf{Metric}} & \textsf{\textbf{Region}}\\
\midrule
\endhead
\multicolumn{6}{@{}l@{}}{\cellcolor{Mist}\textsf{\textbf{\textcolor{Deep}{Neuro / mental health}}}}\\[0.6mm]
\href{https://doi.org/10.1001/jamanetworkopen.2026.25015}{Dales} & Wildfire-sourced PM$_{2.5}$ and ED visits for migraine and other primary headache & \textit{JAMA Netw Open} & Case-crossover & PM$_{2.5}$ (wildfire) & Canada\\
\href{https://doi.org/10.1093/ageing/afag215}{Glover} & Umbrella meta-review of urban environmental exposures and cognitive health & \textit{Age Ageing} & Umbrella review / meta-analysis & PM$_{2.5}$ & Global\\
\href{https://doi.org/10.1016/j.nsa.2026.107019}{Goossens} & In utero air pollution exposure and pre-/postnatal brain development & \textit{Neurosci Appl} & Narrative review & PM$_{2.5}$ / PM$_{10}$ / PM$_{\mathrm{coarse}}$ & Global\\
\href{https://doi.org/10.1016/j.isci.2026.116617}{Ma} & PM$_{2.5}$-induced neurotoxicity and depression: exposure risks and mechanisms & \textit{iScience} & Narrative review & PM$_{2.5}$ + components & Global\\
\href{https://doi.org/10.3760/cma.j.cn112150-20250804-00752}{Pan} & Maternal PM exposure and offspring neurodevelopment: effects and mechanisms & \textit{Zhonghua Yu Fang Yi Xue Za Zhi} & Narrative review & PM (bulk) & China\\
\addlinespace[0.9mm]
\multicolumn{6}{@{}l@{}}{\cellcolor{Mist}\textsf{\textbf{\textcolor{Teal}{Cardiovascular \& metabolic}}}}\\[0.6mm]
\href{https://doi.org/10.3760/cma.j.cn112338-20251120-00843}{Huang} & PM$_{2.5}$ and dyslipidaemia in rural Chinese children; dietary effect modification & \textit{Zhonghua Liu Xing Bing Xue Za Zhi} & Cross-sectional (multistage) & PM$_{2.5}$ & China\\
\href{https://doi.org/10.32394/pe/218966}{Kazmierska} & Air pollution and cardiovascular disease: review of current data & \textit{Przegl Epidemiol} & Narrative review & PM$_{2.5}$ / PM$_{10}$ & Global\\
\href{https://doi.org/10.1007/s00125-026-06814-2}{Xi} & Air pollution, built environment and temperature vs. prediabetes phenotypes & \textit{Diabetologia} & Prospective cohort & PM$_{2.5}$ & Germany\\
\href{https://doi.org/10.1161/JAHA.125.046917}{Yang} & Short-term PM$_{2.5}$ and hospitalisation/ED use in atrial fibrillation & \textit{J Am Heart Assoc} & Case-crossover & PM$_{2.5}$ (wildfire-influenced) & USA\\
\addlinespace[0.9mm]
\multicolumn{6}{@{}l@{}}{\cellcolor{Mist}\textsf{\textbf{\textcolor{Sky}{Reproductive \& developmental}}}}\\[0.6mm]
\href{https://doi.org/10.1016/j.jes.2025.10.043}{Chen} & Prenatal air pollution, fetal growth trajectory phenotypes and child neurodevelopment & \textit{J Environ Sci (China)} & Prospective cohort & PM$_{2.5}$ / PM$_{10}$ & China\\
\href{https://doi.org/10.1016/j.ecoenv.2026.120566}{Huang} & Joint prenatal bisphenol + air-pollutant mixture exposure and preschool IQ & \textit{Ecotoxicol Environ Saf} & Prospective cohort (mixtures) & PM$_{2.5}$ / PM$_{10}$ & China\\
\href{https://doi.org/10.1007/s40572-026-00552-8}{Peterson} & Ambient air pollution, PAHs and pubertal development: critical review & \textit{Curr Environ Health Rep} & Critical review & PM$_{2.5}$ components & Global\\
\href{https://doi.org/10.1016/j.envres.2026.125313}{Xie} & Long-term PM$_{2.5}$ and timing of pubertal development in boys and girls & \textit{Environ Res} & Prospective cohort & PM$_{2.5}$ & China\\
\addlinespace[0.9mm]
\multicolumn{6}{@{}l@{}}{\cellcolor{Mist}\textsf{\textbf{\textcolor{Sage}{Respiratory \& allergic}}}}\\[0.6mm]
\href{https://doi.org/10.1007/s00484-026-03289-z}{Bayar Muluk} & Climate change and otolaryngology: airway disease impacts and emission reduction & \textit{Int J Biometeorol} & Narrative review & PM$_{2.5}$ & Global\\
\href{https://doi.org/10.4103/mgr.MEDGASRES-D-25-00164}{Liao} & Short-term air pollution, respiratory symptoms and lung function in asthma & \textit{Med Gas Res} & Panel study & PM$_{2.5}$ / PM$_{10}$ & China\\
\href{https://doi.org/10.3390/epidemiologia7040097}{Sanchez de Prada} & Ambient PM and asthma-coded admissions in children 0-3 y, Spain & \textit{Epidemiologia} & Ecological panel & PM$_{2.5}$ / PM$_{10}$ & Spain\\
\href{https://doi.org/10.1016/j.jhazmat.2026.143076}{Yang} & Allergenic vs. toxic risk of airborne seed hairs compared with ambient PM$_{2.5}$ & \textit{J Hazard Mater} & Field measurement & PM$_{2.5}$ + bioaerosol & China\\
\addlinespace[0.9mm]
\multicolumn{6}{@{}l@{}}{\cellcolor{Mist}\textsf{\textbf{\textcolor{Violet}{Mechanistic toxicology}}}}\\[0.6mm]
\href{https://doi.org/10.1016/j.jes.2025.12.074}{Ding} & PM$_{2.5}$ promotes vascular calcification via miR-27b-5p/HES1 signalling & \textit{J Environ Sci (China)} & In vitro + organoid + murine & PM$_{2.5}$ & China\\
\href{https://doi.org/10.1007/s00360-026-01698-5}{Maraschi} & Osmoregulatory disruption and metal bioaccumulation from settleable atmospheric PM & \textit{J Comp Physiol B} & Controlled exposure (ecotox) & Settleable APM & Brazil\\
\href{https://doi.org/10.1016/j.jes.2025.11.018}{Wu} & Macrophage NLRP3 inflammasome in real-ambient PM-induced hepatic glucose dysregulation & \textit{J Environ Sci (China)} & Real-ambient murine + co-culture & Real-ambient PM & China\\
\addlinespace[0.9mm]
\multicolumn{6}{@{}l@{}}{\cellcolor{Mist}\textsf{\textbf{\textcolor{Clay}{Occupational \& indoor}}}}\\[0.6mm]
\href{https://doi.org/10.3389/fmedt.2026.1796828}{Chen} & Nested suction instrument to cut intraoperative surgical-smoke PM$_{2.5}$ & \textit{Front Med Technol} & Bench / device experiment & PM$_{2.5}$ (surgical smoke) & China\\
\href{https://doi.org/10.1093/annweh/wxag055}{Cole} & Determinants of respirable crystalline silica exposure in Australian tunnelling & \textit{Ann Work Expo Health} & Observational exposure study & Respirable dust / RCS & Australia\\
\href{https://doi.org/10.1016/j.jhazmat.2026.143081}{Ding} & Stage-resolved hair Zn-isotope dynamics as a firefighting combustion-exposure tracer & \textit{J Hazard Mater} & Repeated-measures biomarker & Combustion-derived PM & China\\
\href{https://doi.org/10.3390/jox16040119}{Ramsden} & Xenobiotic hazards in aircraft cabin air & \textit{J Xenobiot} & Review + risk economics & Ultrafine particles & Global\\
\addlinespace[0.9mm]
\multicolumn{6}{@{}l@{}}{\cellcolor{Mist}\textsf{\textbf{\textcolor{Amber}{Sensing, forecasting \& instrumentation}}}}\\[0.6mm]
\href{https://doi.org/10.1016/j.jenvman.2026.130515}{Shi} & Adaptive spatiotemporal graph transformer for multi-site multi-step PM$_{2.5}$ forecasting & \textit{J Environ Manage} & Deep-learning model & PM$_{2.5}$ & China / India\\
\href{https://doi.org/10.3390/bios16070380}{Tsuchiya} & Seasonal and indoor-outdoor characteristics of fluorescent aerosol particles in offices & \textit{Biosensors} & Real-time sensor campaign & Bioaerosol / coarse & Japan\\
\addlinespace[0.9mm]
\multicolumn{6}{@{}l@{}}{\cellcolor{Mist}\textsf{\textbf{\textcolor{Amber}{Exposure assessment \& modelling}}}}\\[0.6mm]
\href{https://doi.org/10.1029/2025GH001621}{Solorzano-Araque} & Multiscale exposure modelling: LOTOS-EUROS + mobility surveys + particle toxicity & \textit{GeoHealth} & CTM + mobility integration & PM$_{2.5}$ / PM$_{10}$ & Colombia\\
\addlinespace[0.9mm]
\multicolumn{6}{@{}l@{}}{\cellcolor{Mist}\textsf{\textbf{\textcolor{Coral}{Burden, policy \& mitigation}}}}\\[0.6mm]
\href{https://doi.org/10.1002/hsr2.72886}{Karamzad} & Ischaemic heart disease burden in MENA 1990-2021; PM attributable fraction 28.2\% & \textit{Health Sci Rep} & GBD secondary analysis & Ambient PM & MENA\\
\href{https://doi.org/10.1371/journal.pone.0347989}{Rees} & HEPA filtration to reduce respiratory infection in care homes & \textit{PLoS One} & RCT process evaluation & Filtered PM & UK\\
\href{https://doi.org/10.1016/j.jes.2026.01.009}{Shan} & Policy-driven emission reductions and post-pandemic rebound, high-resolution inventories & \textit{J Environ Sci (China)} & Emission inventory & PM emissions & China\\
\addlinespace[0.9mm]
\multicolumn{6}{@{}l@{}}{\cellcolor{Mist}\textsf{\textbf{\textcolor{Slate}{Other clinical endpoints}}}}\\[0.6mm]
\href{https://doi.org/10.3390/dj14070440}{Lupita} & PM$_{2.5}$, preventive-care capacity and caries/edentulism burden across the EU & \textit{Dent J} & Ecological panel & PM$_{2.5}$ & EU-27\\
\href{https://doi.org/10.1016/j.jdin.2026.06.008}{Park} & Long-term PM exposure and risk of androgenetic alopecia & \textit{JAAD Int} & Nationwide cohort & PM$_{2.5}$ / PM$_{10}$ & South Korea\\
\href{https://doi.org/10.3171/2026.3.PEDS2615}{Sun} & Environmental exposure and area deprivation vs. paediatric hydrocephalus outcomes & \textit{J Neurosurg Pediatr} & Registry cohort + geospatial & PM$_{2.5}$ (satellite-derived) & USA\\
\addlinespace[0.9mm]
\bottomrule
\end{longtable}}


\vspace{4mm}
\begin{keybox}
{\sffamily\bfseries\footnotesize\color{Deep}Provenance and method}\\[1.2mm]
{\fontsize{7.2}{8.8}\selectfont
Entry window \textbf{31 July 2026}, contiguous with the 30 July issue --- no gap.
Sources: \textbf{PubMed} (E-utilities and connector) queried on the \texttt{[Date -
Entry]} field along two independent axes, health and instrumentation; \textbf{Europe PMC}
\texttt{CREATION\_DATE:[2026-07-31 TO 2026-07-31]}; \textbf{arXiv} relevance query with a
local recency screen (none in window); \textbf{OpenAlex} returned HTTP 429 and remains
unusable for date-windowed harvesting; \textbf{ClinicalTrials.gov} returned one record
updated in the window --- NCT05160948, personal-monitor result feedback as a behavioural
intervention, observational, so no endpoint analysis was run --- accumulated to
\texttt{state/trials.json} for the weekly rollup. \textbf{11 records in scope.}
\\[1.0mm]
\textbf{Scheduling correction.} This issue was originally built with a two-day window
(31 July--1 August) and published as the 1 August issue. That was wrong: papers entering
PubMed on 1 August belong to the 1 August run. The window was re-cut to 31 July on the
PubMed \texttt{entrez} entry date of every record. \textbf{Four records with entry date
2026-08-01} --- Zhang (PM$_{2.5}$ components and cognition), Daoud (suicide mortality),
Bello (prenatal PM$_{2.5}$, Quebec) and Pai (incense burning and SGA) --- were
\emph{removed} from this issue, dropped from \texttt{state/seen.json} so they will be
re-found, and stashed in \texttt{state/deferred\_2026-08-01.json} for the 1 August run.
Every summary retained here is unchanged from its original text.
\\[1.0mm]
Two records reached this issue through \textbf{Europe PMC} rather than PubMed and carry
earlier PubMed entry dates (Knapova, 8 July; Balu, 21 July). They are dated to this issue
because 31 July is when the Europe PMC creation-date window first surfaced them; a
PubMed-only pipeline would never have recovered them. This is stated rather than hidden
because it means the issue's window is not uniform across sources.
\\[1.0mm]
\textbf{Known defect: zero sensing records for the second consecutive window.} PubMed does
not index \textit{Atmos.\ Meas.\ Tech.} or \textit{Atmos.\ Chem.\ Phys.}; the Consensus
connector has no entry-date filter; OpenAlex date filtering is paywalled. The July
monthly rollup quantified the cost --- 23 sensing records for the month against 11 across
all five daily issues. The fix remains a Crossref by-journal feed keyed on
\texttt{created} for AMT (1867-8548) and ACP (1680-7324).
\\[1.0mm]
Subtopic, design, particle-metric and geography labels were assigned from title,
abstract, keywords and MeSH terms; assignment is single-label by dominant contribution
and is therefore a judgement, not a controlled vocabulary. Effect estimates in the forest
plot are transcribed verbatim from abstracts and use heterogeneous exposure contrasts ---
see the figure caption. Every DOI was verified against the PMID's own PubMed record and
resolved before publication. Content derived from PubMed and Europe PMC metadata;
abstracts remain the copyright of their respective publishers.\par}
\end{keybox}

\end{document}
